\documentclass[aspectratio=169,10pt]{beamer}

\usetheme{Madrid}
\usecolortheme{dolphin}
\setbeamertemplate{navigation symbols}{}
\setbeamertemplate{footline}[frame number]

\usepackage[T1]{fontenc}
\usepackage[utf8]{inputenc}
\usepackage{lmodern}
\usepackage{booktabs}
\usepackage{array}
\usepackage{listings}
\usepackage{xcolor}

\definecolor{leanblue}{RGB}{24,78,119}
\definecolor{leangreen}{RGB}{42,157,143}
\definecolor{leanorange}{RGB}{230,126,34}
\definecolor{leanred}{RGB}{190,70,70}
\definecolor{leanbg}{RGB}{248,249,250}

\setbeamercolor{structure}{fg=leanblue}
\setbeamercolor{frametitle}{fg=white,bg=leanblue}
\setbeamercolor{title}{fg=white,bg=leanblue}
\setbeamercolor{block title}{fg=white,bg=leanblue}
\setbeamercolor{block body}{bg=leanbg}

\lstdefinelanguage{Lean}{
  morekeywords={def,theorem,inductive,structure,where,namespace,end,import,match,with,by,do,let,if,then,else,return,pure,some,none,Prop,Type,Option,List,Array,mutual,partial,private,syntax,macro_rules,example,axiom},
  sensitive=true,
  morecomment=[l]{--},
  morecomment=[s]{/-}{-/},
  morestring=[b]"
}

\lstset{
  language=Lean,
  basicstyle=\ttfamily\scriptsize,
  keywordstyle=\color{leanblue}\bfseries,
  commentstyle=\color{gray},
  stringstyle=\color{leangreen},
  backgroundcolor=\color{leanbg},
  frame=single,
  breaklines=true,
  columns=fullflexible,
  keepspaces=true,
  showstringspaces=false
}

\newcommand{\file}[1]{\texttt{\detokenize{#1}}}
\newcommand{\lean}[1]{\texttt{\detokenize{#1}}}
\newcommand{\statusproved}{\textcolor{leangreen}{proved}}
\newcommand{\statusexec}{\textcolor{leanblue}{executable}}
\newcommand{\statusadmit}{\textcolor{leanred}{admitted}}
\newcommand{\statusfuture}{\textcolor{leanorange}{future work}}

\title[CLean Advisor Update]{CLean: Lean Semantics, PTX Ingestion, and the Saxpy Proof Path}
\subtitle{Architecture, methodology, implementation, and remaining proof obligations}
\author{Riyaza}
\institute{Technical presentation for PhD advisor}
\date{\today}

\begin{document}

% 1
\begin{frame}
  \titlepage
\end{frame}

% 2
\begin{frame}{Deck Roadmap}
  \begin{enumerate}
    \item High-level research idea and design goals.
    \item Core implementation: types, state, instructions, typing.
    \item Semantics: executable helper layer and relational small-step layer.
    \item Proof methodology: computation lemmas, frame lemmas, determinism, decomposition.
    \item PTX front end: AST, parser, checked lowering, bridge theorem surface.
    \item Saxpy case study: PTX input, state construction, correctness pipeline, current blockers.
  \end{enumerate}
\end{frame}

% 3
\begin{frame}{One-Sentence Thesis}
  CLean is a Lean 4 framework for giving a proof-facing, executable semantics to a supported subset of PTX-like GPU kernels, then using that semantics to prove concrete and eventually parametric kernel correctness.

  \vspace{1em}
  \begin{block}{Current repo scope}
    The repo already has a semantic state model, an IR, executable evaluation functions, relational small-step rules, PTX parsing and checked lowering, proof infrastructure, and case studies for toy kernels, saxpy, and matmul.
  \end{block}
\end{frame}

% 4
\begin{frame}{Why This Is Useful}
  \begin{itemize}
    \item GPU kernels are usually validated by testing, compiler inspection, and informal reasoning.
    \item PTX is low-level enough to expose memory spaces, thread IDs, predicates, branches, barriers, and instruction effects.
    \item Lean lets us state exactly what a parsed-and-lowered kernel computes.
    \item Executable semantics let small concrete cases run by \lean{native_decide}; relational semantics support proof statements over arbitrary traces.
  \end{itemize}
\end{frame}

% 5
\begin{frame}{Design Principle}
  The project avoids building a ``PTX text interpreter'' as the trusted core.

  \vspace{0.5em}
  Instead:
  \begin{itemize}
    \item parse raw PTX text into a syntax AST,
    \item checked-lower that AST into a smaller semantic IR,
    \item execute the IR against a concrete state model,
    \item prove properties of the IR semantics,
    \item expose parse/lower soundness as the trusted boundary.
  \end{itemize}
\end{frame}

% 6
\begin{frame}{Repo Topology}
  \begin{tabular}{>{\ttfamily}p{0.32\linewidth}p{0.58\linewidth}}
    \toprule
    CLean/Core & scalar types, values, addresses, IR, state, typing \\
    CLean/Semantics & helper evaluator, small-step relation, executable runner \\
    CLean/Proof & frame lemmas, computation lemmas, loop API, determinism, decomposition \\
    CLean/PTX & PTX AST, parser, lowering, checked bridge \\
    CLean/Examples & semantic smoke tests, parser/lowering tests, saxpy, matmul \\
    \bottomrule
  \end{tabular}
\end{frame}

% 7
\begin{frame}{Status Legend}
  \begin{tabular}{>{\bfseries}p{0.22\linewidth}p{0.68\linewidth}}
    \toprule
    \statusproved & Lean proof has no local \lean{sorry}. \\
    \statusexec & executable function or concrete example checked by \lean{native_decide}. \\
    \statusadmit & theorem statement exists, but body contains \lean{sorry}; this is an explicit proof obligation. \\
    \statusfuture & concept is architected but not implemented or not yet proof-complete. \\
    \bottomrule
  \end{tabular}

  \vspace{1em}
  The deck uses these statuses to separate implemented machinery from admitted proof surface.
\end{frame}

% 8
\begin{frame}{Main Architectural Pipeline}
  \begin{center}
    PTX text
    \quad $\rightarrow$ \quad
    \file{PTX.Ast}
    \quad $\rightarrow$ \quad
    checked lowering
    \quad $\rightarrow$ \quad
    CLean IR
    \quad $\rightarrow$ \quad
    small-step semantics
    \quad $\rightarrow$ \quad
    kernel postcondition
  \end{center}

  \vspace{1em}
  The core proofs are intended to live after lowering, where the instruction set is smaller and semantic cases are explicit.
\end{frame}

% 9
\begin{frame}{Methodological Split}
  \begin{itemize}
    \item \textbf{Executable layer:} functions returning \lean{Option State}, used for simulation, tracing, and concrete regression tests.
    \item \textbf{Relational layer:} inductive propositions \lean{StepInstr}, \lean{StepBlock}, \lean{StepWarp}, \lean{StepMachine}.
    \item \textbf{Bridge layer:} soundness theorems from executable steps to relational steps.
    \item \textbf{Correctness layer:} \lean{PartialCorrect}, \lean{TotalCorrect}, loop specs, and kernel-specific postconditions.
  \end{itemize}
\end{frame}

% 10
\begin{frame}{Implemented Surface Area}
  \begin{itemize}
    \item Scalar integer, predicate, bitwise, address conversion, load/store, branch, barrier, and simple float operations are represented.
    \item Warp-level operations, atomics, and MMA are present in the IR shape, but most are not semantically executed yet.
    \item PTX parser supports declarations, labels, predicated branches, params, shared declarations, memory operations, arithmetic, comparisons, and unsupported-instruction reporting.
    \item Examples validate end-to-end parser, lowering, and execution on small kernels.
  \end{itemize}
\end{frame}

% 11
\begin{frame}{Current Proof Surface}
  \begin{itemize}
    \item \statusproved: executable-step soundness into \lean{StepMachine}.
    \item \statusproved: runN reachability and stuck-stability theorems.
    \item \statusproved: single-warp determinism bridge, modulo preservation lemmas used beneath it.
    \item \statusadmit: several frame and preservation lemmas in \file{Proof/Lemmas.lean}.
    \item \statusadmit: parser/lowering soundness theorem surfaces in \file{Proof/Bridge.lean}.
    \item \statusadmit: parametric saxpy symbolic execution lemma.
  \end{itemize}
\end{frame}

% 12
\begin{frame}{Advisor-Level Takeaway}
  The project has moved from design sketch to an executable, typed, PTX-ingesting semantic core.

  \vspace{0.5em}
  The remaining work is not a vague ``prove everything'' task. It decomposes into:
  \begin{itemize}
    \item finish frame/preservation lemmas,
    \item complete lane-disjointness/decomposition,
    \item do straight-line symbolic execution for saxpy,
    \item package parser/lowering soundness,
    \item then scale to loop kernels like matmul.
  \end{itemize}
\end{frame}

\section{Core Model}

% 13
\begin{frame}{Core Files}
  \begin{tabular}{>{\ttfamily}p{0.32\linewidth}p{0.58\linewidth}}
    \toprule
    Core/Types.lean & scalar types, values, address spaces, special registers \\
    Core/IR.lean & semantic instruction vocabulary \\
    Core/State.lean & machine state, lane/warp/CTA/global memory \\
    Core/Typing.lean & computable type signatures and memory preconditions \\
    \bottomrule
  \end{tabular}
\end{frame}

% 14
\begin{frame}{Scalar Type Universe}
  \begin{itemize}
    \item \lean{ScalarTy} includes predicates, unsigned/signed integers, bit-vectors, half/bfloat/f32/f64.
    \item The semantic layer uses scalar tags directly rather than encoding everything as raw bitvectors.
    \item This keeps type-directed memory encoding and instruction support explicit.
    \item Unsupported or not-yet-modeled cases can fail by returning \lean{none} rather than silently inventing behavior.
  \end{itemize}
\end{frame}

% 15
\begin{frame}[fragile]{Scalar Types in Lean}
\begin{lstlisting}
inductive ScalarTy where
  | pred
  | u8 | u16 | u32 | u64
  | s8 | s16 | s32 | s64
  | b8 | b16 | b32 | b64
  | f16 | bf16 | f32 | f64
  deriving Repr, DecidableEq, Inhabited
\end{lstlisting}

  This is intentionally close to PTX scalar suffixes: \lean{.u32}, \lean{.s32}, \lean{.b64}, \lean{.f32}, etc.
\end{frame}

% 16
\begin{frame}{Address Spaces}
  \begin{itemize}
    \item PTX state spaces are not erased into one integer address space.
    \item \lean{AddrSpace} tracks \lean{global}, \lean{shared}, \lean{local}, \lean{param}, \lean{const}, and \lean{generic}.
    \item Concrete \lean{Addr} constructors carry the scope needed for shared and local memory.
    \item Generic addresses carry a tag and offset, letting \lean{cvta} and \lean{isspacep} have semantic content.
  \end{itemize}
\end{frame}

% 17
\begin{frame}[fragile]{Address Model}
\begin{lstlisting}
inductive Addr where
  | global (offset : Nat)
  | shared (cta : CTAId) (offset : Nat)
  | local (cta : CTAId) (warp : WarpId) (lane : LaneId) (offset : Nat)
  | param (offset : Nat)
  | const (offset : Nat)
  | generic (space : AddrSpace) (offset : Nat)
\end{lstlisting}

  This is one of the most important proof choices: state-space identity is part of the address.
\end{frame}

% 18
\begin{frame}{Value Model}
  \begin{itemize}
    \item \lean{Value} mirrors scalar types and includes semantic addresses.
    \item Signed integer values are stored as Lean \lean{Int}, with normalization helpers for fixed-width behavior.
    \item Floating values are represented using Lean \lean{Float}; this is convenient but not bit-exact PTX f32 semantics.
    \item Tensor fragments are represented by \lean{Value.frag}, but MMA execution is not yet implemented.
  \end{itemize}
\end{frame}

% 19
\begin{frame}[fragile]{Selected Value Constructors}
\begin{lstlisting}
inductive Value where
  | pred (b : Bool)
  | u32 (x : UInt32)
  | u64 (x : UInt64)
  | s32 (x : Int)
  | s64 (x : Int)
  | f32 (x : Float)
  | f64 (x : Float)
  | gaddr (space : AddrSpace) (offset : Nat)
  | frag (ty : FragTy) (payload : Array Value)
\end{lstlisting}
\end{frame}

% 20
\begin{frame}{Thread and Grid Identity}
  \begin{itemize}
    \item \lean{CTAId} and \lean{WarpId} are natural numbers.
    \item \lean{LaneId := Fin 32}, so lane identity is bounded at the type level.
    \item \lean{GridCtx} stores \lean{gridDim} and \lean{blockDim}.
    \item Helpers compute \lean{tid.x}, \lean{ctaid.x}, \lean{ntid.x}, and related special registers from this context.
  \end{itemize}
\end{frame}

% 21
\begin{frame}{Memory Model}
  \begin{itemize}
    \item Every memory space is modeled as a byte map: \lean{Std.HashMap Nat UInt8}.
    \item Scalar loads/stores encode and decode little-endian byte sequences.
    \item Typed memory access requires codec support, byte width, alignment, and address-space agreement.
    \item This gives precise failure points for unsupported or ill-typed accesses.
  \end{itemize}
\end{frame}

% 22
\begin{frame}[fragile]{Top-Level State}
\begin{lstlisting}
structure State where
  kernelEnv : KernelEnv := {}
  global : GlobalMem := {}
  const : ConstMem := {}
  param : ParamMem := {}
  ctas : Std.HashMap CTAId CTAState := {}
  atomics : AtomicState := {}
\end{lstlisting}

  The top-level state separates immutable-ish launch environment from mutable memory and CTA state.
\end{frame}

% 23
\begin{frame}[fragile]{CTA, Warp, Lane}
\begin{lstlisting}
structure LaneState where
  regs : Std.HashMap RegName Value := {}
  preds : Std.HashMap PredName Bool := {}
  localMem : LocalMem := {}
  pc : PC := ("entry", 0)
  status : LaneStatus := .running

structure WarpState where
  lanes : Array LaneState := Array.replicate 32 {}
  activeMask : UInt32 := 0xFFFFFFFF
  exitedMask : UInt32 := 0
\end{lstlisting}
\end{frame}

% 24
\begin{frame}{Well-Formedness Checks}
  \begin{itemize}
    \item \lean{WarpState.wf?}: lane array size is exactly 32.
    \item \lean{CTAState.wf?}: every warp in the CTA is well formed.
    \item \lean{KernelEnv.wf?}: block map keys agree with block labels.
    \item \lean{State.wf?}: kernel environment is well formed and every CTA is well formed.
  \end{itemize}
\end{frame}

% 25
\begin{frame}{Why Boolean and Prop Versions Exist}
  \begin{itemize}
    \item Functions like \lean{stepInstr?} need executable booleans.
    \item Relational rules and proofs need propositions.
    \item The project provides bridge theorems like \lean{wf_iff_bool}.
    \item This pattern repeats for \lean{RunnablePc}, \lean{ParticipatingRunnable}, and \lean{lockstepRunnable}.
  \end{itemize}
\end{frame}

% 26
\begin{frame}{Instruction Vocabulary}
  \begin{itemize}
    \item Pure expression layer: \lean{RValue}.
    \item Guarded instruction layer: \lean{GInstr}.
    \item Block layer: array of guarded instructions plus a terminator.
    \item Effects include register/predicate writes, load/store, address conversion, space tests, barriers, warp ops, atomics, and MMA.
  \end{itemize}
\end{frame}

% 27
\begin{frame}[fragile]{RValue}
\begin{lstlisting}
inductive RValue where
  | imm (v : Value)
  | reg (r : RegName)
  | pred (p : PredName)
  | special (s : SpecialReg)
  | unop (op : ScalarUnaryOp) (a : RValue)
  | binop (op : ScalarBinaryOp) (a b : RValue)
  | triop (op : ScalarTernaryOp) (a b c : RValue)
\end{lstlisting}

  This is the pure expression language used by both instructions and branch conditions.
\end{frame}

% 28
\begin{frame}[fragile]{Instruction Shape}
\begin{lstlisting}
inductive Instr where
  | assignReg (dst : RegName) (rhs : RValue)
  | assignPred (dst : PredName) (cmp : CmpExpr)
  | assignPredValue (dst : PredName) (rhs : RValue)
  | load (dst : RegName) (src : TypedAddr)
  | store (dst : TypedAddr) (value : RValue)
  | cvta (dst : RegName) (space : AddrSpace) (src : RValue)
  | isspacep (dst : PredName) (space : AddrSpace) (src : RValue)
  | barrierCTA (barrierId : Nat)
  | warp (op : WarpOp)
  | atomic ...
  | mma ...
\end{lstlisting}
\end{frame}

% 29
\begin{frame}{Guarded Instructions}
  \begin{itemize}
    \item \lean{GInstr} pairs an optional guard with an instruction.
    \item A guard is a predicate register plus an optional negation bit.
    \item Participation is computed lane-by-lane against runnable lanes at the current PC.
    \item Predicated-off lanes at barriers are explicitly advanced rather than blocked.
  \end{itemize}
\end{frame}

% 30
\begin{frame}{Terminators and Blocks}
  \begin{itemize}
    \item \lean{Terminator.br}: unconditional branch to a block label.
    \item \lean{Terminator.cbr}: conditional branch with true and false labels.
    \item \lean{Terminator.terminate}: marks current runnable lanes as terminated.
    \item A \lean{Block} is a label, body array, and terminator.
  \end{itemize}
\end{frame}

% 31
\begin{frame}{Typing Layer}
  \begin{itemize}
    \item \file{Core/Typing.lean} is the repo's local type checker for the semantic IR and lowering phase.
    \item It defines byte widths and alignments.
    \item It computes signatures for unary, binary, ternary, and comparison operations.
    \item It defines typed memory access preconditions.
  \end{itemize}
\end{frame}

% 32
\begin{frame}[fragile]{Typed Access Preconditions}
\begin{lstlisting}
def typedAccessPreconditions? (space : AddrSpace)
    (ty : ScalarTy) (addr : Addr) : Bool :=
  scalarCodecSupported? ty &&
  (byteWidth? ty).isSome &&
  aligned? ty addr &&
  addrSpaceMatches? space addr
\end{lstlisting}

  This check is used by both \lean{readMem?} and \lean{writeMem?}.
\end{frame}

% 33
\begin{frame}{Scalar Compatibility}
  \begin{itemize}
    \item Checked lowering accepts some PTX-compatible aliases: for example \lean{s32}, \lean{u32}, and \lean{b32}.
    \item \lean{scalarCompatible?} captures these expected/actual relationships.
    \item \lean{coerceOperandTo?} inserts conversions when the semantic value needs an explicit shape.
    \item 64-bit pointer-shaped values are deliberately kept unwrapped in some cases because runtime \lean{cvta} may produce generic addresses.
  \end{itemize}
\end{frame}

% 34
\begin{frame}{Supported Operation Signatures}
  \begin{itemize}
    \item \lean{unarySig?}: moves, negation, abs, bitnot, conversion.
    \item \lean{binarySig?}: add, sub, mul, wide multiply, bitwise logic, shifts, min/max.
    \item \lean{ternarySig?}: \lean{mad}, \lean{fma}, \lean{selp}.
    \item \lean{cmpSig?}: same-type integer/float comparisons returning predicate.
  \end{itemize}
\end{frame}

% 35
\begin{frame}{Core Model Summary}
  \begin{itemize}
    \item The state model is concrete enough to execute kernels.
    \item The IR is small enough to reason about directly.
    \item The typing layer rejects unsupported or inconsistent PTX before semantics.
    \item The major abstraction boundary is already in place: PTX syntax lowers into semantic effects.
  \end{itemize}
\end{frame}

\section{Semantics}

% 36
\begin{frame}{Semantics Files}
  \begin{tabular}{>{\ttfamily}p{0.35\linewidth}p{0.55\linewidth}}
    \toprule
    Semantics/Helpers.lean & executable evaluator and state transformers \\
    Semantics/SmallStep.lean & relational small-step rules \\
    Semantics/Execution.lean & reachability, step driver, runN, traceN \\
    \bottomrule
  \end{tabular}
\end{frame}

% 37
\begin{frame}{Executable Helper Layer}
  \begin{itemize}
    \item Most semantic content lives in functions ending with \lean{?}.
    \item Failure is represented by \lean{Option.none}.
    \item This keeps unsupported cases, bad lookups, invalid alignment, and divergent branch conditions explicit.
    \item The same functions are used for concrete execution and for proof rules.
  \end{itemize}
\end{frame}

% 38
\begin{frame}{Runnable Lanes}
  \begin{itemize}
    \item \lean{laneIsRunnable} checks lane status and active mask bit.
    \item \lean{runnableLaneIds} filters all 32 lanes.
    \item \lean{currentRunnablePc?} takes the first runnable lane's PC.
    \item \lean{lockstepRunnable?} requires all runnable lanes to have the same PC.
  \end{itemize}
\end{frame}

% 39
\begin{frame}{Participation}
  \begin{itemize}
    \item \lean{participatingRunnableLaneIds?} starts from runnable lanes.
    \item It filters to lanes at the current PC.
    \item It evaluates the optional guard per lane.
    \item It returns the lane list that actually participates in the instruction effect.
  \end{itemize}
\end{frame}

% 40
\begin{frame}{Special Registers}
  \begin{itemize}
    \item \lean{evalSpecial} maps PTX special registers to values.
    \item Thread IDs come from block-local linear lane identity plus \lean{blockDim}.
    \item CTA IDs are decoded from a linear \lean{CTAId} and \lean{gridDim}.
    \item \lean{ntid} and \lean{nctaid} come directly from launch context.
  \end{itemize}
\end{frame}

% 41
\begin{frame}{Expression Evaluation}
  \begin{itemize}
    \item \lean{evalRValue?} is mutually recursive with scalar operator evaluation.
    \item Register and predicate reads are lane-local.
    \item Special registers are computed from \lean{GridCtx}, CTA, warp, and lane.
    \item Pure operators implement arithmetic, logical, conversion, and address arithmetic behavior.
  \end{itemize}
\end{frame}

% 42
\begin{frame}[fragile]{RValue Evaluation Skeleton}
\begin{lstlisting}
partial def evalRValue? st cta warp lane : RValue -> Option Value
  | .imm v => some v
  | .reg r => do
      let laneState <- st.getLane? cta warp lane
      readReg laneState r
  | .pred p => do
      let laneState <- st.getLane? cta warp lane
      let b <- readPred laneState p
      pure (.pred b)
  | .special s => some <| evalSpecial st.kernelEnv.gridCtx cta warp lane s
  | .binop op a b => ...
\end{lstlisting}
\end{frame}

% 43
\begin{frame}{Memory Read and Write}
  \begin{itemize}
    \item \lean{readMem?} checks typed access preconditions.
    \item It finds the correct byte map for the address space.
    \item It reads the required byte width and decodes the scalar.
    \item \lean{writeMem?} performs the dual encode/write/update operation.
  \end{itemize}
\end{frame}

% 44
\begin{frame}{Address Resolution}
  \begin{itemize}
    \item \lean{resolveAddr?} evaluates a \lean{TypedAddr}'s address expression.
    \item Integer-like offsets lower into concrete state-space addresses.
    \item Generic addresses resolve only when their tag matches the requested space, except \lean{.generic} keeps the tag.
    \item Shared and local addresses add CTA, warp, and lane scope.
  \end{itemize}
\end{frame}

% 45
\begin{frame}{Address Conversion}
  \begin{itemize}
    \item \lean{evalCvta?} maps integer or already-generic pointer-shaped values to \lean{Value.gaddr}.
    \item \lean{evalIsspacep?} tests the tag of a generic address.
    \item This supports kernels that load pointers from parameter memory before global memory access.
    \item Saxpy relies on this path: parameter pointers are loaded, converted to global addresses, offset, and dereferenced.
  \end{itemize}
\end{frame}

% 46
\begin{frame}{Instruction Step Function}
  \begin{itemize}
    \item \lean{stepInstr?} fetches warp state and checks lockstep.
    \item It computes participating lanes under the optional guard.
    \item For register/predicate operations, it applies a lane-local update over participants.
    \item For stores, it sequentializes participating lanes in lane-list order.
    \item It then advances runnable PCs, except for barriers which have special handling.
  \end{itemize}
\end{frame}

% 47
\begin{frame}[fragile]{stepInstr? Shape}
\begin{lstlisting}
def stepInstr? (st : State) (cta : CTAId) (warp : WarpId)
    (gi : GInstr) : Option State := do
  let warpState <- st.getWarp? cta warp
  if !lockstepRunnable? warpState then none else
    let participants <- participatingRunnableLaneIds? warpState gi.guard?
    match gi.instr with
    | .barrierCTA barrierId => stepBarrierCTA? st cta warp barrierId participants
    | _ => do
        let st <- match gi.instr with
          | .assignReg dst rhs => ...
          | .load dst src => ...
          | .store dst value => ...
          | _ => none
        advanceRunnablePcs? st cta warp
\end{lstlisting}
\end{frame}

% 48
\begin{frame}{Barrier Semantics}
  \begin{itemize}
    \item \lean{stepBarrierCTA?} models a CTA barrier instance keyed by barrier id.
    \item Predicated-off lanes skip the barrier and advance.
    \item Participating lanes record arrival and become \lean{blockedBarrier}.
    \item Once arrivals reach expected count, all arrived lanes are released, advanced, and marked running.
  \end{itemize}
\end{frame}

% 49
\begin{frame}{Terminator Semantics}
  \begin{itemize}
    \item \lean{br}: sets PC of current lanes to target block index 0.
    \item \lean{terminate}: marks current lanes terminated.
    \item \lean{cbr}: evaluates condition for all current lanes and requires a uniform destination.
    \item Non-uniform branch conditions currently fail by returning \lean{none}.
  \end{itemize}
\end{frame}

% 50
\begin{frame}{Relational Small-Step Layer}
  \begin{itemize}
    \item \lean{StepInstr} wraps \lean{stepInstr?} with preconditions.
    \item \lean{StepBlock} chooses between body instruction and block terminator.
    \item \lean{StepWarp} packages a block step for a CTA/warp.
    \item \lean{StepMachine} is the global one-step relation.
  \end{itemize}
\end{frame}

% 51
\begin{frame}[fragile]{StepInstr Rule}
\begin{lstlisting}
inductive StepInstr : State -> CTAId -> WarpId -> GInstr -> State -> Prop where
  | mk :
      State.wf st ->
      st.getWarp? cta warp = some warpState ->
      WarpState.wf warpState ->
      Helpers.lockstepRunnable warpState ->
      Helpers.ParticipatingRunnable warpState gi.guard? participants ->
      Helpers.stepInstr? st cta warp gi = some st' ->
      StepInstr st cta warp gi st'
\end{lstlisting}
\end{frame}

% 52
\begin{frame}{Executable Driver}
  \begin{itemize}
    \item \lean{currentInstrStep?}: attempt the current body instruction.
    \item \lean{currentTermStep?}: attempt the terminator when body index is past the body.
    \item \lean{stepAt?}: instruction first, then terminator.
    \item \lean{step?}: default scheduler for CTA 0, warp 0.
  \end{itemize}
\end{frame}

% 53
\begin{frame}{runN and traceN}
  \begin{itemize}
    \item \lean{runN fuel st} executes up to \lean{fuel} successful \lean{step?} transitions.
    \item If \lean{step?} returns \lean{none}, \lean{runN} stays at the current state.
    \item \lean{runN?} fails if the requested number of steps cannot be taken.
    \item \lean{traceN} records intermediate states for debugging and smoke tests.
  \end{itemize}
\end{frame}

% 54
\begin{frame}{Reachability}
  \begin{itemize}
    \item \lean{Reaches} is the reflexive-transitive closure of \lean{StepMachine}.
    \item \lean{StepMachine.runN_reaches}: every executable run is a relational reachability trace.
    \item \lean{StepMachine.runN?_sound}: successful exact-fuel execution is also relationally reachable.
    \item This is the base connection between executable testing and proof-facing semantics.
  \end{itemize}
\end{frame}

% 55
\begin{frame}{Semantics Summary}
  The semantics now has two synchronized views:
  \begin{itemize}
    \item a computable interpreter for concrete kernels;
    \item an inductive step relation for theorem statements.
  \end{itemize}

  \vspace{0.5em}
  The key engineering challenge is keeping these two views aligned without duplicating semantic definitions.
\end{frame}

\section{Proof Methodology}

% 56
\begin{frame}{Proof Files}
  \begin{tabular}{>{\ttfamily}p{0.36\linewidth}p{0.54\linewidth}}
    \toprule
    Proof/Lemmas.lean & frame, memory, wf, executable-to-relational helpers \\
    Proof/InstrCompute.lean & per-instruction computation over lane lists \\
    Proof/Loops.lean & correctness definitions and counted-loop specs \\
    Proof/Determinism.lean & single-warp determinism bridge \\
    Proof/IsSingleWarpPres.lean & preservation of single-warp support \\
    Proof/LaneDecomposition.lean & lane-local view and lifting architecture \\
    Proof/Bridge.lean & parser/lowering soundness theorem surface \\
    \bottomrule
  \end{tabular}
\end{frame}

% 57
\begin{frame}{Correctness Vocabulary}
  \begin{itemize}
    \item \lean{MachineFinal st}: no \lean{StepMachine} successor exists.
    \item \lean{TerminatesAt init final}: \lean{Reaches init final} and \lean{MachineFinal final}.
    \item \lean{PartialCorrect init post}: every terminal final state satisfies \lean{post}.
    \item \lean{TotalCorrect init post}: some terminal final state exists and satisfies \lean{post}.
  \end{itemize}
\end{frame}

% 58
\begin{frame}[fragile]{Correctness Definitions}
\begin{lstlisting}
def MachineFinal (st : State) : Prop :=
  forall st', not (StepMachine st st')

def TerminatesAt (init final : State) : Prop :=
  Reaches init final /\ MachineFinal final

def PartialCorrect (init : State) (post : State -> Prop) : Prop :=
  forall final, TerminatesAt init final -> post final
\end{lstlisting}
\end{frame}

% 59
\begin{frame}{Computation Lemma Pattern}
  \begin{itemize}
    \item First prove structural unfoldings of executable loops.
    \item Then prove one-lane or list-of-lanes effects.
    \item Then package these effects as instruction-specific computation lemmas.
    \item This is especially important because \lean{stepInstr?} uses imperative-style \lean{for} loops inside \lean{Option}.
  \end{itemize}
\end{frame}

% 60
\begin{frame}{applyToLaneIds? Foundation}
  \begin{itemize}
    \item \lean{applyToLaneIds?_nil}: empty lane list returns the input state.
    \item \lean{applyToLaneIds?_cons}: exposes one state update plus recursion.
    \item \lean{applyToLaneIds?_singleton}: common one-lane simplification.
    \item \lean{applyToLaneIds?_lane_in}: for nodup lists, participating lane gets the computed state.
  \end{itemize}
\end{frame}

% 61
\begin{frame}{Frame Lemma Pattern}
  \begin{itemize}
    \item Register writes preserve predicate files, local memory, PC, and status.
    \item Predicate writes preserve register files, local memory, PC, and status.
    \item Lane updates preserve other lanes.
    \item CTA/warp/lane state updates preserve top-level global/const/param fields.
  \end{itemize}
\end{frame}

% 62
\begin{frame}{Memory Lemmas}
  \begin{itemize}
    \item \lean{natToBytesLE_length}: encoded byte list has requested width.
    \item \lean{readBytes?_writeBytes_same}: writing bytes then reading same range returns those bytes.
    \item \lean{decode_encode_*}: scalar roundtrip lemmas for supported types.
    \item \lean{readMem_writeMem_same_global_*}: typed global write/read roundtrips.
  \end{itemize}
\end{frame}

% 63
\begin{frame}{State Preservation Lemmas}
  \begin{itemize}
    \item Many instructions should not mutate global/param/const memory.
    \item Store should preserve lane register and predicate files.
    \item Barriers should preserve data memory and register/predicate files, but mutate barrier and lane status/PC.
    \item These facts become reusable frame facts for kernel proofs.
  \end{itemize}
\end{frame}

% 64
\begin{frame}{Current Lemma Status}
  \begin{itemize}
    \item \statusproved: basic register/predicate read-after-write and non-alias frame lemmas.
    \item \statusproved: byte encoding, reading, and global memory read-after-write for key scalar types.
    \item \statusproved: non-store instruction top-memory preservation for several cases.
    \item \statusadmit: barrier frame lemmas, several reg/pred preservation lemmas, generic wf preservation, and kernelEnv preservation.
  \end{itemize}
\end{frame}

% 65
\begin{frame}{Executable-to-Relational Soundness}
  \begin{itemize}
    \item \lean{body_of_stepInstr?}: if helper instruction step succeeds under preconditions, produce \lean{StepMachine}.
    \item \lean{term_of_stepTerminator?}: same for terminators.
    \item \lean{body_of_currentInstrStep?_computed}: executable current instruction step is sound.
    \item \lean{term_of_currentTermStep?_computed}: executable current terminator step is sound.
  \end{itemize}
\end{frame}

% 66
\begin{frame}{Automation Layer}
  \begin{itemize}
    \item \lean{cstep}: discharges concrete \lean{StepMachine} goals by using current-step soundness and \lean{native_decide}.
    \item \lean{crun}: uses \lean{runN_reaches}.
    \item \lean{cpost}: calls \lean{native_decide} on boolean postconditions.
    \item \lean{cfinish}: tries one of the above.
  \end{itemize}
\end{frame}

% 67
\begin{frame}{Why Determinism Is Scoped}
  \begin{itemize}
    \item General GPU scheduling is intentionally not solved yet.
    \item The current executable driver steps CTA 0, warp 0.
    \item Many examples, including saxpy, use a single populated warp.
    \item The project proves that, under \lean{IsSingleWarp}, any relational step is captured by \lean{step?}.
  \end{itemize}
\end{frame}

% 68
\begin{frame}[fragile]{IsSingleWarp}
\begin{lstlisting}
def IsSingleWarp (st : State) : Prop :=
  forall cta warp ws,
    st.getWarp? cta warp = some ws -> cta = 0 /\ warp = 0
\end{lstlisting}

  This predicate says the only warp that exists in the state is CTA 0, warp 0.
\end{frame}

% 69
\begin{frame}{Single-Warp Determinism Theorem}
  \begin{itemize}
    \item \lean{step?_of_StepMachine}: under \lean{IsSingleWarp}, any \lean{StepMachine st st'} implies \lean{step? st = some st'}.
    \item \lean{reaches_imp_runN_of_IsSingleWarp}: any reachable final state is some \lean{runN fuel init}.
    \item \lean{terminal_eq_runN}: terminal states reduce to executable runs.
    \item This is the core bridge used by saxpy partial correctness.
  \end{itemize}
\end{frame}

% 70
\begin{frame}{Stuck Stability}
  \begin{itemize}
    \item \lean{step?_eq_none_of_MachineFinal}: final relational state is stuck for the executable scheduler.
    \item \lean{runN_eq_of_step?_none}: once stuck, extra fuel does nothing.
    \item \lean{runN_add}: executable runs compose over fuel.
    \item \lean{runN_eq_of_both_MachineFinal}: two final \lean{runN} endpoints from the same init are equal.
  \end{itemize}
\end{frame}

% 71
\begin{frame}{Single-Warp Preservation}
  \begin{itemize}
    \item \file{Proof/IsSingleWarpPres.lean} proves updates do not introduce new warp keys.
    \item The main invariant is \lean{WarpSupportEq}: two states agree on which CTA/warp lookups succeed.
    \item Primitive updates such as \lean{setLane}, \lean{writeMem?}, and \lean{advanceRunnablePcs?} preserve support.
    \item \lean{step?_preserves_IsSingleWarp} is the headline preservation theorem.
  \end{itemize}
\end{frame}

% 72
\begin{frame}{Single-Warp Preservation Status}
  \begin{itemize}
    \item \statusproved: most primitive support-preservation lemmas.
    \item \statusproved: instruction, terminator, and \lean{step?} preservation structure.
    \item \statusadmit: \lean{stepBarrierCTA?_preserves_warp_support} is still a \lean{sorry}.
    \item Consequence: the high-level preservation theorem currently depends on that admitted barrier lemma.
  \end{itemize}
\end{frame}

% 73
\begin{frame}{Loop Proof API}
  \begin{itemize}
    \item \lean{CountedLoopSpec} abstracts an invariant, guard, body relation, variant, and postcondition.
    \item \lean{partial_correct}: invariant preservation plus exit rule implies postcondition.
    \item \lean{total_correct}: decreasing natural variant gives termination and postcondition.
    \item This is mostly for future loop-heavy kernels like matmul; saxpy is straight-line after the early-exit branch.
  \end{itemize}
\end{frame}

% 74
\begin{frame}{Lane Decomposition Goal}
  \begin{itemize}
    \item GPU warp semantics updates multiple lanes at once.
    \item Saxpy correctness is naturally per element, hence per lane.
    \item Need to show each lane's local view evolves as expected under full-warp execution.
    \item Need to combine per-lane facts into a global memory postcondition.
  \end{itemize}
\end{frame}

% 75
\begin{frame}{Lane View}
  \begin{itemize}
    \item \lean{LaneView} projects one lane's registers, predicates, local memory, PC, and status.
    \item \lean{LaneLocal st lane} extracts the view at CTA 0, warp 0.
    \item \lean{instrWriteSet?} computes addresses written by a lane for store/atomic instructions.
    \item \lean{DisjointLaneWrites} states pairwise write-set disjointness across active lanes.
  \end{itemize}
\end{frame}

% 76
\begin{frame}{Lane Decomposition Status}
  \begin{itemize}
    \item \statusproved: definitions for write sets, disjointness, lane view, and lift are implemented.
    \item \statusproved: current commutation statements are weak/existence-shaped and do not yet prove saxpy.
    \item \statusfuture: strengthen \lean{lanes_independent_runN} to a real per-lane run theorem.
    \item \statusfuture: connect lane-disjoint stores to global vector postconditions.
  \end{itemize}
\end{frame}

% 77
\begin{frame}{Proof Methodology Summary}
  \begin{itemize}
    \item Keep executable semantics as the single semantic source.
    \item Prove executable steps are sound relational steps.
    \item Use frame lemmas to localize instruction effects.
    \item Use determinism for single-warp kernels to reduce arbitrary terminal traces to \lean{runN}.
    \item Use lane decomposition to lift per-lane symbolic execution to full-warp vector correctness.
  \end{itemize}
\end{frame}

\section{PTX Front End}

% 78
\begin{frame}{PTX Files}
  \begin{tabular}{>{\ttfamily}p{0.33\linewidth}p{0.57\linewidth}}
    \toprule
    PTX/Ast.lean & syntax AST independent of semantic IR \\
    PTX/Parser.lean & parser combinators with source-position errors \\
    PTX/Lowering.lean & checked lowering to CLean IR \\
    PTX/Bridge.lean & parse-plus-lower API and error type \\
    Proof/Bridge.lean & intended soundness theorem surface \\
    \bottomrule
  \end{tabular}
\end{frame}

% 79
\begin{frame}{PTX AST Role}
  \begin{itemize}
    \item It preserves PTX syntactic distinctions before semantic decisions.
    \item Operands include registers, predicates, symbols, addresses, immediates, and special registers.
    \item Instructions include supported operations and an \lean{unsupported} fallback.
    \item Kernels store declarations, params, shared declarations, and parsed blocks.
  \end{itemize}
\end{frame}

% 80
\begin{frame}[fragile]{PTX Instruction AST}
\begin{lstlisting}
inductive Instr where
  | mov (ty : ScalarTy) (dst : RegName) (src : Operand)
  | binop (op : ScalarBinaryOp) (ty : ScalarTy) (dst : RegName) (lhs rhs : Operand)
  | setp (op : CmpOp) (ty : ScalarTy) (dst : PredName) (lhs rhs : Operand)
  | ld (space : AddrSpace) (ty : ScalarTy) (dst : RegName) (addr : Operand)
  | st (space : AddrSpace) (ty : ScalarTy) (addr value : Operand)
  | cvta (space : AddrSpace) (dst : RegName) (src : Operand)
  | barSync (barrierId : Nat)
  | unsupported (opcode : String) (modifiers : Array String) (operands : Array Operand)
\end{lstlisting}
\end{frame}

% 81
\begin{frame}{Parser Design}
  \begin{itemize}
    \item Implemented with Lean's \lean{Std.Internal.Parsec.String}.
    \item Tracks line and column on parse failure.
    \item Handles whitespace, line comments, block comments, and pragmas.
    \item Parses module directives, declarations, labels, instructions, and kernels.
  \end{itemize}
\end{frame}

% 82
\begin{frame}{Parser Instruction Coverage}
  \begin{itemize}
    \item Arithmetic: \lean{add}, \lean{sub}, \lean{mul.lo}, \lean{mul.wide.s32}, \lean{mad.lo}, \lean{fma.rn}.
    \item Predicate ops: \lean{setp}, \lean{or.pred}, \lean{and.pred}, \lean{xor.pred}.
    \item Memory: \lean{ld.<space>.<ty>}, \lean{st.<space>.<ty>}.
    \item Address: \lean{cvta.to.<space>}, \lean{isspacep.<space>}.
    \item Control: labels, \lean{bra}, guarded \lean{@p bra}, \lean{exit}/\lean{ret}.
  \end{itemize}
\end{frame}

% 83
\begin{frame}{Parser Block Construction}
  \begin{itemize}
    \item Statements are parsed as instructions, terminators, guarded branches, or no-ops.
    \item Guarded branches split the current statement stream into a block ending in \lean{cbra}.
    \item Fallthrough blocks receive generated labels.
    \item Implicit fallthroughs are linked to the next parsed block or to exit.
  \end{itemize}
\end{frame}

% 84
\begin{frame}{Lowering Strategy}
  \begin{itemize}
    \item Unchecked lowering exists for direct syntax-to-IR mapping.
    \item Checked lowering is the important path.
    \item Checked lowering threads a \lean{Typing.TypeEnv}.
    \item Each instruction checks operands, supported types, guards, labels, params, and shared symbols.
  \end{itemize}
\end{frame}

% 85
\begin{frame}{Lowering Error Model}
  \begin{itemize}
    \item \lean{unknownReg}, \lean{unknownPred}, \lean{unknownParam}, \lean{unknownShared}.
    \item \lean{unknownSymbol}, \lean{unknownBlock}.
    \item \lean{typeMismatch}, \lean{unsupportedType}, \lean{unsupportedOp}, \lean{unsupportedModifier}.
    \item This is crucial for making unsupported PTX fail loudly at the boundary.
  \end{itemize}
\end{frame}

% 86
\begin{frame}{Checked Lowering Examples}
  \begin{itemize}
    \item \lean{mov.u32} lowers to \lean{assignReg}.
    \item \lean{add.u32} lowers to \lean{assignReg} with \lean{RValue.binop add}.
    \item \lean{setp.ge.s32} lowers to \lean{assignPred}.
    \item \lean{ld.global.s32} lowers to \lean{load} with a typed global address.
    \item \lean{st.global.s32} lowers to \lean{store}.
  \end{itemize}
\end{frame}

% 87
\begin{frame}[fragile]{Lowering Skeleton}
\begin{lstlisting}
def lowerInstrChecked? (env : Typing.TypeEnv) :
    Instr -> LowerM (CLean.Instr x Typing.TypeEnv)
  | .mov ty dst src => do
      expectOperandType env ty src
      let src <- lowerOperandCheckedAs? env ty src
      pure (.assignReg dst src, { env with regs := env.regs.insert dst ty })
  | .ld space ty dst addr => do
      ensureCodecType ty
      let addrTy <- addressOperandType? env space addr
      ...
\end{lstlisting}

  Note: the actual file uses Lean product syntax; this slide uses \lean{x} visually for readability.
\end{frame}

% 88
\begin{frame}{Kernel Environment Lowering}
  \begin{itemize}
    \item Parameters are laid out with \lean{lowerParamsChecked?}.
    \item Shared declarations are laid out with \lean{lowerSharedsChecked?}.
    \item Blocks are lowered into a \lean{HashMap BlockLabel Block}.
    \item Entry label is checked against the block-label map.
  \end{itemize}
\end{frame}

% 89
\begin{frame}{Parser/Lowering Bridge}
  \begin{itemize}
    \item \lean{BridgeError} distinguishes parse errors from lowering errors.
    \item \lean{parseAndLowerKernel?} parses one kernel and checked-lowers it.
    \item \lean{parseAndLowerModule?} checks every kernel in a parsed module.
    \item Boolean wrappers power examples: \lean{parseAndLowerKernelOk?}, \lean{parseAndLowerModuleOk?}.
  \end{itemize}
\end{frame}

% 90
\begin{frame}{Proof Bridge Status}
  \begin{itemize}
    \item \file{Proof/Bridge.lean} defines \lean{LoweredKernelWf} and \lean{LoweredModuleWf}.
    \item \statusadmit: \lean{lowerKernelEnvChecked_wf}.
    \item \statusadmit: \lean{parseAndLowerKernel?_sound}.
    \item \statusadmit: \lean{parseAndLowerModule?_sound}.
    \item These are clean theorem surfaces, but not yet fully proved.
  \end{itemize}
\end{frame}

\section{Examples and Regression Tests}

% 91
\begin{frame}{Example Strategy}
  \begin{itemize}
    \item Start with tiny semantic IR kernels: assign, copy, barrier.
    \item Then test manually built PTX AST lowering.
    \item Then parse PTX text and lower it.
    \item Finally execute parsed/lowered kernels with \lean{runN} and boolean postconditions.
  \end{itemize}
\end{frame}

% 92
\begin{frame}{Semantic Examples}
  \begin{itemize}
    \item Toy assign: one lane writes \lean{r1 = 7} and terminates.
    \item Toy copy: load global \lean{u32}, store to another global offset, terminate.
    \item Toy barrier: barrier releases under expected count 1, then assignment and termination.
    \item These examples use \lean{cstep}, \lean{runN}, and \lean{native_decide}.
  \end{itemize}
\end{frame}

% 93
\begin{frame}{Lowering Examples}
  \begin{itemize}
    \item Direct PTX AST instructions are lowered and compared by reflexivity.
    \item Checked lowering rejects use-before-declare.
    \item Parsed/lowered PTX kernels run under the semantic driver.
    \item Coverage includes assign, add, copy, param copy, barrier, and shared memory.
  \end{itemize}
\end{frame}

% 94
\begin{frame}{Executable Validation Philosophy}
  \begin{itemize}
    \item Concrete examples do not prove parametric correctness.
    \item They are still valuable because they exercise the actual parser, lowering, and semantics.
    \item \lean{native_decide} catches many integration bugs quickly.
    \item The proof architecture then explains what remains beyond testing.
  \end{itemize}
\end{frame}

\section{Saxpy Case Study}

% 95
\begin{frame}{Saxpy Target}
  The saxpy-style integer kernel computes:
  \[
    r[i] = x[i] * \alpha + y[i]
  \]
  for active lanes whose thread index is below \lean{n}.

  \vspace{0.5em}
  In this repo, the data values are \lean{s32}; the postcondition uses 32-bit signed wrapping through \lean{s32Wrap}.
\end{frame}

% 96
\begin{frame}{Saxpy PTX Shape}
  \begin{itemize}
    \item Parameters: \lean{n}, \lean{alpha}, \lean{x} pointer, \lean{y} pointer, \lean{r} pointer.
    \item Loads parameter values from \lean{.param}.
    \item Computes \lean{threadIdx = ctaid.x * ntid.x + tid.x}.
    \item Branches to return if \lean{threadIdx >= n}.
    \item Otherwise loads \lean{x[i]} and \lean{y[i]}, computes \lean{x[i] * alpha + y[i]}, stores \lean{r[i]}.
  \end{itemize}
\end{frame}

% 97
\begin{frame}[fragile]{Saxpy PTX Core}
\begin{lstlisting}
ld.param.u32 %r2, [saxpyKernel_param_0];
ld.param.s32 %r6, [saxpyKernel_param_1];
ld.param.u64 %rd1, [saxpyKernel_param_2];
...
mad.lo.s32 %r1, %r3, %r4, %r5;
setp.ge.s32 %p1, %r1, %r2;
@%p1 bra $L__BB0_2;
...
ld.global.s32 %r7, [%rd6];
ld.global.s32 %r8, [%rd8];
mad.lo.s32 %r9, %r7, %r6, %r8;
st.global.s32 [%rd10], %r9;
\end{lstlisting}
\end{frame}

% 98
\begin{frame}{Saxpy Parsing and Lowering}
  \begin{itemize}
    \item \lean{saxpyKernelText} stores the PTX source.
    \item \lean{PTX.parseAndLowerKernelOk? saxpyKernelText = true} is checked by \lean{native_decide}.
    \item \lean{PTX.lowerKernelSupported? saxpyKernel = true} is also checked.
    \item This validates that the current parser and checked-lowering subset covers this kernel.
  \end{itemize}
\end{frame}

% 99
\begin{frame}{Saxpy Initial State}
  \begin{itemize}
    \item \lean{saxpyXBase = 0}, \lean{saxpyYBase = 128}, \lean{saxpyRBase = 256}.
    \item Parameter memory stores \lean{n}, \lean{alpha}, and three base pointers.
    \item Global memory stores input vectors \lean{xs} and \lean{ys}.
    \item One CTA and one warp are populated; active mask is \lean{activeMaskPrefix n}.
  \end{itemize}
\end{frame}

% 100
\begin{frame}[fragile]{saxpyStateFor}
\begin{lstlisting}
def saxpyStateFor (n : Nat) (alpha : Int) (xs ys : List Int) : State :=
  { kernelEnv := PTX.lowerKernelEnvCheckedD saxpyKernel
    global := { bytes := saxpyGlobalBytesFor xs ys }
    param := { bytes := saxpyParamBytesFor n alpha }
    ctas := ({} : Std.HashMap CTAId CTAState).insert 0
      { warps := ({} : Std.HashMap WarpId WarpState).insert 0
          (saxpyWarpFor n) } }
\end{lstlisting}
\end{frame}

% 101
\begin{frame}{Saxpy Postcondition}
  \begin{itemize}
    \item \lean{saxpyExpectedAt alpha xs ys i} computes wrapped \lean{xs[i] * alpha + ys[i]}.
    \item \lean{saxpyPost base n alpha xs ys st} requires every index below \lean{n} to match expected global memory.
    \item It is built from \lean{vectorS32Post}.
    \item This is the parametric final property for saxpy.
  \end{itemize}
\end{frame}

% 102
\begin{frame}[fragile]{Saxpy Spec Helpers}
\begin{lstlisting}
def saxpyExpectedAt (alpha : Int) (xs ys : List Int) (i : Nat) : Int :=
  s32Wrap (listIntGetD xs i * alpha + listIntGetD ys i)

def saxpyPost (base n : Nat) (alpha : Int) (xs ys : List Int)
    (st : State) : Prop :=
  vectorS32Post base n (saxpyExpectedAt alpha xs ys) st
\end{lstlisting}
\end{frame}

% 103
\begin{frame}{Concrete Saxpy Status}
  \begin{itemize}
    \item File header documents an axiom-clean concrete run for \lean{n = 4}.
    \item Expected output: \lean{[12, 24, 36, 48]} for \lean{alpha = 2}, \lean{xs = [1,2,3,4]}, \lean{ys = [10,20,30,40]}.
    \item The current visible file focuses on the general theorem path and the admitted symbolic execution lemma.
    \item Running \file{CLean/Examples/Saxpy.lean} currently succeeds with a warning that \lean{saxpy_runN_satisfies_post} uses \lean{sorry}.
  \end{itemize}
\end{frame}

% 104
\begin{frame}{Saxpy Single-Warp Lemma}
  \begin{itemize}
    \item \lean{saxpyStateFor_isSingleWarp} is proved.
    \item It unfolds the state construction and proves any successful warp lookup must be CTA 0, warp 0.
    \item This is the entry point into the single-warp determinism machinery.
    \item It is specific to the state constructor, not an axiom about the kernel.
  \end{itemize}
\end{frame}

% 105
\begin{frame}{Saxpy Proof Reduction}
  The general \lean{saxpy_partial_correct} proof has a meaningful skeleton:

  \begin{enumerate}
    \item Initial state is \lean{IsSingleWarp}.
    \item \lean{step?} preserves \lean{IsSingleWarp}.
    \item Any terminal relational final equals \lean{runN K init}.
    \item A symbolic execution theorem supplies a terminal \lean{runN J init} satisfying the post.
    \item Stuck-stability equates the two terminal \lean{runN} endpoints.
  \end{enumerate}
\end{frame}

% 106
\begin{frame}{Remaining Saxpy Bottleneck}
  The key remaining theorem is:

  \vspace{0.5em}
  \lean{saxpy_runN_satisfies_post}

  \vspace{0.5em}
  It must show that for every supported launch shape \lean{n <= 32} with matching input lengths, there exists fuel \lean{K} such that \lean{runN K init} is final and satisfies \lean{saxpyPost}.
\end{frame}

% 107
\begin{frame}{What the Saxpy RunN Theorem Requires}
  \begin{itemize}
    \item A straight-line symbolic execution of the saxpy body.
    \item Per-lane reasoning for thread index, branch predicate, pointer offsets, loads, multiply-add, and store.
    \item Read-after-write memory lemmas for each lane's output slot.
    \item Lane-disjointness proof: lane \lean{i} writes only output address \lean{rBase + i * 4}.
    \item Lift from per-lane stores to vector postcondition.
  \end{itemize}
\end{frame}

% 108
\begin{frame}{Saxpy Early-Exit Case}
  \begin{itemize}
    \item For lanes with \lean{i >= n}, \lean{setp.ge.s32} sets the branch predicate.
    \item The guarded branch takes them to return.
    \item These lanes should not write output.
    \item The postcondition only quantifies \lean{i < n}, so inactive/out-of-range behavior is framed away.
  \end{itemize}
\end{frame}

% 109
\begin{frame}{Saxpy Active-Lane Case}
  \begin{itemize}
    \item For lanes with \lean{i < n}, the branch predicate is false.
    \item The lane converts parameter pointers to global addresses.
    \item It computes byte offset \lean{i * 4}.
    \item It reads \lean{x[i]} and \lean{y[i]}.
    \item It writes \lean{s32Wrap (x[i] * alpha + y[i])} to \lean{r[i]}.
  \end{itemize}
\end{frame}

% 110
\begin{frame}{Saxpy Memory Argument}
  \begin{itemize}
    \item Input and output regions are separated by base constants 0, 128, and 256.
    \item For \lean{n <= 32}, each vector occupies at most 128 bytes.
    \item Output slots \lean{rBase + i * 4} are pairwise distinct.
    \item The write/read lemmas already cover global \lean{s32} stores, but the vector-wide frame/lift is still needed.
  \end{itemize}
\end{frame}

% 111
\begin{frame}{What Is Left Before saxpy Is Honest}
  \begin{itemize}
    \item Fill remaining \file{Proof/Lemmas.lean} frame and preservation holes.
    \item Strengthen \file{Proof/LaneDecomposition.lean} from shape statements to pointwise lane execution.
    \item Prove the active-lane symbolic execution sequence for saxpy.
    \item Prove lane-disjoint store addresses.
    \item Package the postcondition proof and remove \lean{saxpy_runN_satisfies_post}'s \lean{sorry}.
  \end{itemize}
\end{frame}

% 112
\begin{frame}{Saxpy Progress Report}
  \begin{tabular}{p{0.50\linewidth}p{0.30\linewidth}}
    \toprule
    Component & Status \\
    \midrule
    PTX source parses & \statusexec \\
    Checked lowering accepts kernel & \statusexec \\
    State constructor and memory layout & implemented \\
    Single-warp initial-state lemma & \statusproved \\
    Determinism reduction for partial correctness & \statusproved, with dependencies \\
    Symbolic run theorem & \statusadmit \\
    Total correctness theorem & depends on admitted symbolic run \\
    \bottomrule
  \end{tabular}
\end{frame}

\section{Known Limitations}

% 113
\begin{frame}{Known Semantic Limitations}
  \begin{itemize}
    \item General warp scheduling and multi-CTA interleavings are not solved.
    \item Non-uniform conditional branches currently fail rather than model divergence/reconvergence.
    \item Warp ops, atomics, and MMA are represented but mostly not implemented in \lean{stepInstr?}.
    \item Floating-point semantics use Lean \lean{Float}; this is not bit-exact PTX f32.
  \end{itemize}
\end{frame}

% 114
\begin{frame}{Known Proof Limitations}
  \begin{itemize}
    \item Several frame and preservation theorems are currently admitted.
    \item Parser/lowering soundness is stated but not yet proved.
    \item Lane decomposition is architected but not strong enough for saxpy yet.
    \item Matmul general correctness is intentionally admitted pending loop-invariant proofs.
  \end{itemize}
\end{frame}

% 115
\begin{frame}{Why the Existing Sorries Are Manageable}
  \begin{itemize}
    \item They are named and localized rather than hidden behind broad trusted theorems.
    \item Many have mechanical proof shapes: unfold executable function, split cases, apply update frame lemmas.
    \item The saxpy remaining theorem is kernel-specific and can be decomposed by lane and instruction.
    \item The proof surface makes it clear what \lean{\#print axioms} should report.
  \end{itemize}
\end{frame}

% 116
\begin{frame}{Recommended Next Milestone}
  \begin{enumerate}
    \item Finish \lean{stepInstr?_preserves_kernelEnv} and wf preservation.
    \item Finish barrier frame/support lemmas or restrict saxpy path to no-barrier dependencies where possible.
    \item Prove a strong \lean{applyToLaneIds?} lane-in/lane-not-in theorem sufficient for multi-lane register updates.
    \item Prove saxpy active-lane straight-line computation.
    \item Remove \lean{saxpy_runN_satisfies_post}'s \lean{sorry}.
  \end{enumerate}
\end{frame}

% 117
\begin{frame}{Longer-Term Path}
  \begin{itemize}
    \item Prove parser/lowering well-formedness.
    \item Generalize from single-warp to scheduler-independent multi-CTA reasoning.
    \item Model divergence/reconvergence if needed for real PTX kernels.
    \item Replace float semantics with bit-precise f32/f64 if hardware equivalence is the goal.
    \item Instantiate \lean{CountedLoopSpec} for matmul's unrolled and remainder loops.
  \end{itemize}
\end{frame}

% 118
\begin{frame}{What This Deck Claims}
  \begin{itemize}
    \item CLean already has a functioning semantic pipeline from parsed PTX to executable IR.
    \item Concrete examples are running through the actual semantics.
    \item The theorem architecture for saxpy is in place.
    \item The current blocking work is well-scoped proof engineering, not redesign.
    \item Scaling beyond saxpy will require stronger lane and loop infrastructure.
  \end{itemize}
\end{frame}

% 119
\begin{frame}{Files to Discuss Live}
  \begin{itemize}
    \item \file{CLean/Core/IR.lean}: semantic instruction vocabulary.
    \item \file{CLean/Semantics/Helpers.lean}: executable behavior.
    \item \file{CLean/Proof/Determinism.lean}: single-warp determinism bridge.
    \item \file{CLean/PTX/Parser.lean}: PTX ingestion.
    \item \file{CLean/PTX/Lowering.lean}: checked lowering.
    \item \file{CLean/Examples/Saxpy.lean}: case study and current proof boundary.
  \end{itemize}
\end{frame}

% 120
\begin{frame}{Discussion}
  The central advisor question is now strategic:

  \vspace{0.8em}
  \begin{block}{Next proof target}
    Should the immediate milestone be a fully honest saxpy proof under the current single-warp model, or should the project first strengthen the semantic model for divergence and scheduling before investing in kernel-specific proof?
  \end{block}
\end{frame}

\end{document}
